\documentclass[11pt]{article}
\usepackage[utf8]{inputenc}
\usepackage[T1]{fontenc}
\usepackage[spanish]{babel}
\usepackage{amsmath,amssymb,amsthm,mathtools}
\usepackage{braket}
\usepackage{dsfont}
\usepackage{graphicx}
\usepackage{float}
\usepackage[margin=2.5cm]{geometry}
\usepackage{hyperref}
\hypersetup{hidelinks}
\emergencystretch=2em
\allowdisplaybreaks

\newcommand{\ii}{\mathrm{i}}
\newcommand{\Id}{\mathds{1}}

\title{Entrelazamiento en el modelo HC}
\author{Ricardo Del Rosal Gardu\~no \\ \small Doctorado en F\'isica, BUAP}
\date{Notas de doctorado}

\begin{document}
\maketitle

\tableofcontents
\bigskip

\section{Sistema tripartito y biparticiones}

El modelo de Holstein-Cummings contiene tres grados de libertad: el estado
electr\'onico molecular, el modo fot\'onico de la cavidad y el modo fon\'onico
intramolecular. El espacio de Hilbert completo es, por tanto,
\begin{equation}
  \boxed{
  \mathcal H
  =
  \mathcal H_{\mathrm{el}}
  \otimes
  \mathcal H_{\mathrm{fot}}
  \otimes
  \mathcal H_{\mathrm{fon}}.
  }
\end{equation}

Aunque el sistema es tripartito, primero puede estudiarse como una
bipartici\'on entre el electr\'on y los otros dos grados de libertad:
\begin{equation}
  \mathcal H
  =
  \mathcal H_{\mathrm{el}}
  \otimes
  \left(
  \mathcal H_{\mathrm{fot}}
  \otimes
  \mathcal H_{\mathrm{fon}}
  \right).
\end{equation}
Esta partici\'on permite cuantificar cu\'anto se entrelaza el estado electr\'onico
con el subsistema fot\'on-fon\'on considerado como una sola parte.

Recordando la evoluci\'on obtenida para el hamiltoniano efectivo, el estado en
la representaci\'on desplazada puede reagruparse seg\'un el estado electr\'onico:
\begin{equation}
  \boxed{
  \ket{\tilde\Psi(t)}
  =
  \ket{\chi_g(t)}\ket g
  +
  \ket{\chi_e(t)}\ket e.
  }
\end{equation}
En el sector din\'amico con $n\geq1$ y $m\geq1$, las ramas fot\'on-fon\'on son
\begin{equation}
  \ket{\chi_g(t)}
  =
  \sum_{n=1}^{\infty}
  \sum_{m=1}^{\infty}
  \left[
  p_{nm}^{+}(t)\ket{n,m-1}
  -
  p_{nm}^{-}(t)\ket{n,m}
  \right],
\end{equation}
\begin{equation}
  \ket{\chi_e(t)}
  =
  \sum_{n=1}^{\infty}
  \sum_{m=1}^{\infty}
  \left[
  p_{nm}^{+}(t)\ket{n-1,m-1}
  +
  p_{nm}^{-}(t)\ket{n-1,m}
  \right].
\end{equation}
Estos vectores contienen toda la informaci\'on del campo y de la vibraci\'on
condicionada a encontrar la mol\'ecula en $\ket g$ o en $\ket e$. No se
normalizan por separado. Sus normas son las probabilidades de las ramas
electr\'onicas correspondientes.

\section{Matrices de densidad reducidas}

Como el estado total es puro, su matriz de densidad es
\begin{equation}
  \hat\rho(t)
  =
  \ket{\tilde\Psi(t)}\bra{\tilde\Psi(t)}.
\end{equation}
Sustituyendo la descomposici\'on electr\'onica y desarrollando el producto
exterior se obtiene
\begin{equation}
\begin{aligned}
  \hat\rho(t)
  ={}&
  \ket{\chi_g}\bra{\chi_g}
  \otimes\ket g\bra g
  +
  \ket{\chi_g}\bra{\chi_e}
  \otimes\ket g\bra e
  \\
  &+
  \ket{\chi_e}\bra{\chi_g}
  \otimes\ket e\bra g
  +
  \ket{\chi_e}\bra{\chi_e}
  \otimes\ket e\bra e.
\end{aligned}
\end{equation}

Para estudiar el entrelazamiento entre la parte electr\'onica y el resto se
calcula la matriz de densidad reducida electr\'onica:
\begin{equation}
  \hat\rho_{\mathrm{el}}(t)
  =
  \operatorname{Tr}_{\mathrm{fot,fon}}
  \left[
  \hat\rho(t)
  \right].
\end{equation}
Escribiendo la traza parcial de manera expl\'icita en las bases de Fock de los
modos fot\'onico y fon\'onico,
\begin{equation}
\begin{aligned}
  \hat\rho_{\mathrm{el}}(t)
  =
  \sum_{n,m}
  \Big(&
  \bra n_{\mathrm{fot}}
  \otimes
  \bra m_{\mathrm{fon}}
  \otimes
  \Id_{\mathrm{el}}
  \Big)
  \hat\rho(t)
  \\
  &\times
  \Big(
  \ket n_{\mathrm{fot}}
  \otimes
  \ket m_{\mathrm{fon}}
  \otimes
  \Id_{\mathrm{el}}
  \Big).
\end{aligned}
\end{equation}
La identidad
\begin{equation}
  \operatorname{Tr}(\ket u\bra v)=\braket{v|u}
\end{equation}
convierte cada operador fot\'on-fon\'on en un producto interno. Por tanto,
\begin{equation}
\begin{aligned}
  \hat\rho_{\mathrm{el}}(t)
  ={}&
  \braket{\chi_g|\chi_g}\ket g\bra g
  +
  \braket{\chi_e|\chi_g}\ket g\bra e
  \\
  &+
  \braket{\chi_g|\chi_e}\ket e\bra g
  +
  \braket{\chi_e|\chi_e}\ket e\bra e.
\end{aligned}
\end{equation}
En la base ordenada $\{\ket g,\ket e\}$,
\begin{equation}
  \boxed{
  \hat\rho_{\mathrm{el}}(t)
  =
  \begin{pmatrix}
    \braket{\chi_g|\chi_g} & \braket{\chi_e|\chi_g}
    \\
    \braket{\chi_g|\chi_e} & \braket{\chi_e|\chi_e}
  \end{pmatrix}.
  }
\end{equation}
Los elementos diagonales son las probabilidades electr\'onicas
\begin{equation}
  p_g(t)=\braket{\chi_g|\chi_g},
  \qquad
  p_e(t)=\braket{\chi_e|\chi_e},
  \qquad
  p_g(t)+p_e(t)=1.
\end{equation}

Por otro lado, al trazar sobre el subsistema electr\'onico se obtiene
\begin{equation}
  \hat\rho_{\mathrm{fot-fon}}(t)
  =
  \operatorname{Tr}_{\mathrm{el}}
  \left[
  \hat\rho(t)
  \right].
\end{equation}
Los t\'erminos cruzados desaparecen porque $\braket{g|e}=0$, y queda
\begin{equation}
  \boxed{
  \hat\rho_{\mathrm{fot-fon}}(t)
  =
  \ket{\chi_g(t)}\bra{\chi_g(t)}
  +
  \ket{\chi_e(t)}\bra{\chi_e(t)}.
  }
\end{equation}
Como ambas matrices reducidas provienen de una bipartici\'on de un estado
puro, tienen el mismo espectro no nulo y, en particular, la misma pureza.

\section{Pureza, inversi\'on y coherencia}

La pureza de un subsistema $A$ se define como
\begin{equation}
  P(\hat\rho_A)
  =
  \operatorname{Tr}(\hat\rho_A^2).
\end{equation}
Para un estado reducido puro, $P=1$. En ese caso el subsistema $A$ no est\'a
entrelazado con su complemento. Para un qubit m\'aximamente mezclado, y por
tanto m\'aximamente entrelazado con el resto cuando el estado global es puro,
la pureza toma su valor m\'inimo $P=1/2$.

Definimos la coherencia electr\'onica reducida como
\begin{equation}
  \boxed{
  s(t)=\braket{\chi_g(t)|\chi_e(t)}.
  }
\end{equation}
Entonces
\begin{equation}
  \hat\rho_{\mathrm{el}}
  =
  \begin{pmatrix}
    p_g & s^*\\
    s & p_e
  \end{pmatrix}.
\end{equation}
Al elevar esta matriz al cuadrado,
\begin{equation}
  \hat\rho_{\mathrm{el}}^2
  =
  \begin{pmatrix}
    p_g^2+|s|^2 & p_gs^*+p_es^*
    \\
    p_gs+p_es & p_e^2+|s|^2
  \end{pmatrix}.
\end{equation}
Tomando la traza se obtiene
\begin{equation}
  P_{\mathrm{el}}(t)
  =
  p_g^2(t)+p_e^2(t)+2|s(t)|^2.
\end{equation}

El mismo resultado puede derivarse desde la matriz reducida fot\'on-fon\'on.
En efecto,
\begin{equation}
\begin{aligned}
  \hat\rho_{\mathrm{fot-fon}}^2
  ={}&
  p_g\ket{\chi_g}\bra{\chi_g}
  +s\ket{\chi_g}\bra{\chi_e}
  \\
  &+s^*\ket{\chi_e}\bra{\chi_g}
  +p_e\ket{\chi_e}\bra{\chi_e},
\end{aligned}
\end{equation}
y su traza es
\begin{equation}
  \operatorname{Tr}(\hat\rho_{\mathrm{fot-fon}}^2)
  =
  p_g^2+p_e^2+2|s|^2.
\end{equation}
De manera equivalente, la pureza puede escribirse con las normas de las ramas:
\begin{equation}
  P_{\mathrm{el}}(t)
  =
  \|\chi_g(t)\|^4
  +
  \|\chi_e(t)\|^4
  +
  2\left|\braket{\chi_g(t)|\chi_e(t)}\right|^2.
\end{equation}

La inversi\'on electr\'onica es
\begin{equation}
  W(t)
  =
  p_e(t)-p_g(t).
\end{equation}
Junto con la normalizaci\'on $p_e+p_g=1$, esta ecuaci\'on permite despejar las
poblaciones. Primero,
\begin{equation}
  W+p_g=1-p_g
  \quad\Longrightarrow\quad
  p_g=\frac{1-W}{2}.
\end{equation}
De manera an\'aloga,
\begin{equation}
  p_e-W=1-p_e
  \quad\Longrightarrow\quad
  p_e=\frac{1+W}{2}.
\end{equation}
Por consiguiente,
\begin{equation}
\begin{aligned}
  p_g^2+p_e^2
  &=
  \frac14\left[(1-W)^2+(1+W)^2\right]
  \\
  &=
  \frac14\left(
  1-2W+W^2+1+2W+W^2
  \right)
  \\
  &=
  \frac{1+W^2}{2}.
\end{aligned}
\end{equation}
La pureza electr\'onica general queda
\begin{equation}
  \boxed{
  P_{\mathrm{el}}(t)
  =
  \frac12
  \left[
  1+W^2(t)+4|s(t)|^2
  \right].
  }
\end{equation}
Esta expresi\'on separa la contribuci\'on debida al desbalance de poblaciones,
medido por $W(t)$, de la contribuci\'on debida a la coherencia que permanece
entre las dos ramas electr\'onicas despu\'es de ignorar el campo y la vibraci\'on.

\section{Caso 1: un fon\'on en un bloque polarit\'onico}

Consideremos que el estado permanece dentro de un solo sub-bloque:
\begin{equation}
  \ket{\tilde\Psi(t)}
  =
  u(t)\ket{+,n,m-1}
  +
  v(t)\ket{-,n,m},
\end{equation}
con
\begin{equation}
  |u(t)|^2+|v(t)|^2=1.
\end{equation}
Los estados polarit\'onicos son
\begin{equation}
  \ket{\pm,n}
  =
  \frac{1}{\sqrt2}
  \left(
  \ket{e,n-1}
  \pm
  \ket{g,n}
  \right).
\end{equation}
Al sustituirlos en el estado evolucionado,
\begin{equation}
\begin{aligned}
  \ket{\tilde\Psi(t)}
  ={}&
  \frac{u(t)}{\sqrt2}
  \left(
  \ket{e,n-1}+\ket{g,n}
  \right)
  \ket{m-1}
  \\
  &+
  \frac{v(t)}{\sqrt2}
  \left(
  \ket{e,n-1}-\ket{g,n}
  \right)
  \ket m.
\end{aligned}
\end{equation}
Agrupando los t\'erminos que multiplican a $\ket e$ y a $\ket g$,
\begin{equation}
\begin{aligned}
  \ket{\tilde\Psi(t)}
  ={}&
  \frac{1}{\sqrt2}
  \ket e\ket{n-1}
  \left[
  u(t)\ket{m-1}+v(t)\ket m
  \right]
  \\
  &+
  \frac{1}{\sqrt2}
  \ket g\ket n
  \left[
  u(t)\ket{m-1}-v(t)\ket m
  \right].
\end{aligned}
\end{equation}
Por tanto,
\begin{equation}
  \ket{\chi_e(t)}
  =
  \frac{\ket{n-1}}{\sqrt2}
  \left[
  u(t)\ket{m-1}+v(t)\ket m
  \right],
\end{equation}
\begin{equation}
  \ket{\chi_g(t)}
  =
  \frac{\ket n}{\sqrt2}
  \left[
  u(t)\ket{m-1}-v(t)\ket m
  \right].
\end{equation}
La ortogonalidad de los estados fon\'onicos y la normalizaci\'on de $u$ y $v$
dan
\begin{equation}
  p_e
  =
  \braket{\chi_e|\chi_e}
  =
  \frac12\left(|u|^2+|v|^2\right)
  =
  \frac12,
\end{equation}
\begin{equation}
  p_g
  =
  \braket{\chi_g|\chi_g}
  =
  \frac12.
\end{equation}
La coherencia electr\'onica es nula porque las ramas contienen estados
fot\'onicos ortogonales:
\begin{equation}
  s(t)
  =
  \braket{\chi_g|\chi_e}
  \propto
  \braket{n|n-1}
  =0.
\end{equation}
En consecuencia,
\begin{equation}
  W(t)=p_e-p_g=0,
\end{equation}
y la pureza es
\begin{equation}
  \boxed{
  P_{nm}(t)=\frac12.
  }
\end{equation}
No es indispensable que el estado inicial sea un eigenestado del hamiltoniano
efectivo. Basta con que el estado evolucionado permanezca dentro del mismo
doblete $\{\ket{+,n,m-1},\ket{-,n,m}\}$.

\begin{figure}[H]
  \centering
  \includegraphics[width=0.95\linewidth]{\detokenize{../../notas-doctorado/imgs/Entrelazamiento/pureza-fonon-polariton.png}}
  \caption{Pureza electr\'onica para un fon\'on en un \'unico bloque
  polarit\'onico. La l\'inea permanece en $P_{\mathrm{el}}(t)=1/2$, valor
  m\'inimo de la pureza de un qubit, porque las ramas electr\'onicas tienen
  siempre la misma norma y est\'an acompa\~nadas por estados fot\'onicos
  ortogonales. El electr\'on est\'a m\'aximamente entrelazado con el subsistema
  fot\'on-fon\'on durante toda la evoluci\'on. Las curvas del hamiltoniano HC
  completo y del hamiltoniano efectivo son pr\'acticamente indistinguibles.}
\end{figure}

\section{Caso 2: estado coherente fon\'onico en un bloque polarit\'onico}

Se fija ahora un bloque polarit\'onico $n$, pero el modo vibracional ocupa una
superposici\'on de muchos estados de Fock. El estado evolucionado es
\begin{equation}
\begin{aligned}
  \ket{\tilde\Psi_n(t)}
  ={}&
  p_{n0}^{-}(t)\ket{-,n,0}
  \\
  &+
  \sum_{m=1}^{\infty}
  \left[
  p_{nm}^{+}(t)\ket{+,n,m-1}
  +
  p_{nm}^{-}(t)\ket{-,n,m}
  \right].
\end{aligned}
\end{equation}
El primer t\'ermino es el sector fon\'onico de borde. Expandiendo los
polaritones en la base electr\'onica desnuda,
\begin{equation}
\begin{aligned}
  \ket{\tilde\Psi_n(t)}
  ={}&
  \frac{p_{n0}^{-}(t)}{\sqrt2}
  \left(
  \ket{e,n-1}-\ket{g,n}
  \right)\ket0
  \\
  &+
  \frac{1}{\sqrt2}
  \sum_{m=1}^{\infty}
  p_{nm}^{+}(t)
  \left(
  \ket{e,n-1}+\ket{g,n}
  \right)
  \ket{m-1}
  \\
  &+
  \frac{1}{\sqrt2}
  \sum_{m=1}^{\infty}
  p_{nm}^{-}(t)
  \left(
  \ket{e,n-1}-\ket{g,n}
  \right)
  \ket m.
\end{aligned}
\end{equation}

Reagrupando los t\'erminos que multiplican a $\ket e$,
\begin{equation}
\begin{aligned}
  \ket{\chi_e(t)}
  =
  \frac{\ket{n-1}}{\sqrt2}
  \Bigg[
  &p_{n0}^{-}(t)\ket0
  +
  \sum_{m=1}^{\infty}p_{nm}^{+}(t)\ket{m-1}
  \\
  &+
  \sum_{m=1}^{\infty}p_{nm}^{-}(t)\ket m
  \Bigg].
\end{aligned}
\end{equation}
En la primera suma se introduce $k=m-1$, de modo que $m=k+1$. En la segunda
suma se cambia solamente el nombre del \'indice $m\to k$. Al incorporar el
t\'ermino de borde como $p_{n,0}^{-}(t)$ resulta
\begin{equation}
  \boxed{
  \ket{\chi_e(t)}
  =
  \frac{\ket{n-1}}{\sqrt2}
  \sum_{k=0}^{\infty}
  \left[
  p_{n,k+1}^{+}(t)+p_{n,k}^{-}(t)
  \right]
  \ket k.
  }
\end{equation}
La amplitud asociada al estado fot\'on-fon\'on $\ket{n-1,k}$ es
\begin{equation}
  E_k(t)
  =
  \frac{p_{n,k+1}^{+}(t)+p_{n,k}^{-}(t)}{\sqrt2}.
\end{equation}

Para la rama que acompa\~na a $\ket g$ se obtiene, de manera an\'aloga,
\begin{equation}
\begin{aligned}
  \ket{\chi_g(t)}
  =
  \frac{\ket n}{\sqrt2}
  \Bigg[
  &-p_{n0}^{-}(t)\ket0
  +
  \sum_{m=1}^{\infty}p_{nm}^{+}(t)\ket{m-1}
  \\
  &-
  \sum_{m=1}^{\infty}p_{nm}^{-}(t)\ket m
  \Bigg],
\end{aligned}
\end{equation}
y despu\'es del mismo cambio de \'indices,
\begin{equation}
  \boxed{
  \ket{\chi_g(t)}
  =
  \frac{\ket n}{\sqrt2}
  \sum_{k=0}^{\infty}
  \left[
  p_{n,k+1}^{+}(t)-p_{n,k}^{-}(t)
  \right]
  \ket k.
  }
\end{equation}
La amplitud asociada al estado $\ket{n,k}$ es
\begin{equation}
  G_k(t)
  =
  \frac{p_{n,k+1}^{+}(t)-p_{n,k}^{-}(t)}{\sqrt2}.
\end{equation}

Los estados fot\'onicos que acompa\~nan a las dos ramas electr\'onicas son
ortogonales:
\begin{equation}
  \braket{n|n-1}=0.
\end{equation}
Por ello,
\begin{equation}
  s(t)=\braket{\chi_g|\chi_e}=0.
\end{equation}
La matriz electr\'onica es diagonal y la pureza queda determinada por la
inversi\'on electr\'onica del bloque:
\begin{equation}
  \boxed{
  P_n(t)
  =
  \frac{1+W_n^2(t)}{2}.
  }
\end{equation}

La inversi\'on aparece por interferencia entre amplitudes de dobletes
fon\'onicos adyacentes. En efecto,
\begin{equation}
\begin{aligned}
  W_n(t)
  &=
  p_e(t)-p_g(t)
  \\
  &=
  \sum_{k=0}^{\infty}
  \left(
  |E_k(t)|^2-|G_k(t)|^2
  \right).
\end{aligned}
\end{equation}
Para dos n\'umeros complejos $a$ y $b$ se cumple
\begin{equation}
  \frac12\left(|a+b|^2-|a-b|^2\right)
  =
  2\operatorname{Re}(a^*b).
\end{equation}
Tomando $a=p_{n,k+1}^{+}$ y $b=p_{n,k}^{-}$ se obtiene
\begin{equation}
  \boxed{
  W_n(t)
  =
  2\operatorname{Re}
  \sum_{k=0}^{\infty}
  p_{n,k+1}^{+*}(t)
  p_{n,k}^{-}(t).
  }
\end{equation}

\begin{figure}[H]
  \centering
  \includegraphics[width=0.95\linewidth]{\detokenize{../../notas-doctorado/imgs/Entrelazamiento/pureza-fonon-coherente-polariton.png}}
  \caption{Pureza electr\'onica para un estado fon\'onico coherente dentro de
  un bloque polarit\'onico fijo. Como $s(t)=0$, se cumple exactamente
  $P_n(t)=[1+W_n^2(t)]/2$. Cuando $W_n(t)=0$, las poblaciones $p_e$ y $p_g$
  son iguales y la pureza alcanza $1/2$: el electr\'on est\'a m\'aximamente
  entrelazado con el subsistema fot\'on-fon\'on. Cuando $|W_n(t)|$ se aproxima
  a uno, una poblaci\'on electr\'onica domina, la rama opuesta pierde peso y la
  pureza se acerca a uno. Las oscilaciones r\'apidas provienen de la
  interferencia entre amplitudes fon\'onicas vecinas, mientras que la
  envolvente refleja su desfasamiento y refase colectivos. La evoluci\'on
  obtenida con el hamiltoniano completo y con el efectivo es pr\'acticamente
  id\'entica.}
\end{figure}

\section{Caso 3: estados coherentes fot\'onico y fon\'onico}

Cuando tanto el modo fot\'onico como el fon\'onico comienzan en estados
coherentes aparecen muchos valores de $n$. Las ramas electr\'onicas pueden
compartir estados fot\'on-fon\'on con las mismas etiquetas, de modo que la
coherencia $s(t)$ ya no tiene que anularse.

Reuniendo las amplitudes de las dos ramas en la base fot\'on-fon\'on
$\{\ket{q,k}\}$ se definen
\begin{equation}
  E_{qk}(t)
  =
  \frac{p_{q,k+1}^{+}(t)+p_{q,k}^{-}(t)}{\sqrt2},
  \qquad
  G_{qk}(t)
  =
  \frac{p_{q,k+1}^{+}(t)-p_{q,k}^{-}(t)}{\sqrt2}.
\end{equation}
Aqu\'i $q$ etiqueta el sector fot\'onico despu\'es de reagrupar las ramas en
una base com\'un. Entonces
\begin{equation}
  \ket{\chi_e(t)}
  =
  \sum_{q,k}E_{qk}(t)\ket{q,k},
  \qquad
  \ket{\chi_g(t)}
  =
  \sum_{q,k}G_{qk}(t)\ket{q,k}.
\end{equation}
Las poblaciones electr\'onicas son
\begin{equation}
  p_e(t)
  =
  \braket{\chi_e|\chi_e}
  =
  \sum_{q,k}|E_{qk}(t)|^2,
\end{equation}
\begin{equation}
  p_g(t)
  =
  \braket{\chi_g|\chi_g}
  =
  \sum_{q,k}|G_{qk}(t)|^2.
\end{equation}
La coherencia reducida es
\begin{equation}
  s(t)
  =
  \braket{\chi_g|\chi_e}
  =
  \sum_{q,k}G_{qk}^{*}(t)E_{qk}(t).
\end{equation}
Por tanto, la pureza conserva la forma general
\begin{equation}
  \boxed{
  P_{\alpha,\beta}(t)
  =
  \frac12
  \left[
  1+W^2(t)+4|s(t)|^2
  \right].
  }
\end{equation}
A pesar de que esta expresi\'on es anal\'itica, contiene sumas dobles y las
frecuencias $\Omega_{nm}^{\mathrm{ef}}$ dependen de manera no lineal de $n$
y $m$. En general no existe una suma cerrada y las series deben evaluarse con
un truncamiento num\'erico que abarque las distribuciones coherentes.

\begin{figure}[H]
  \centering
  \includegraphics[width=0.95\linewidth]{\detokenize{../../notas-doctorado/imgs/Entrelazamiento/pureza-coherentes-excitado.png}}
  \caption{Pureza para estados coherentes fot\'onico y fon\'onico con
  $c_e=1$ y $c_g=0$. El estado inicial es un producto y comienza en
  $P_{\mathrm{el}}(0)=1$. La interacci\'on correlaciona al electr\'on con una
  superposici\'on de muchos sectores $(n,m)$, por lo que la pureza desciende y
  desarrolla oscilaciones, colapsos y refases. Los acercamientos a $1/2$
  corresponden a intervalos de entrelazamiento casi m\'aximo entre el electr\'on
  y el subsistema fot\'on-fon\'on. Las soluciones completa y efectiva se
  superponen pr\'acticamente en toda la ventana mostrada.}
\end{figure}

\begin{figure}[H]
  \centering
  \includegraphics[width=0.95\linewidth]{\detokenize{../../notas-doctorado/imgs/Entrelazamiento/pureza-coherentes-simetrico.png}}
  \caption{Pureza para los mismos estados coherentes, ahora con
  $c_e=1/\sqrt2$ y $c_g=-1/\sqrt2$. El estado electr\'onico inicial tambi\'en
  es puro, por lo que $P_{\mathrm{el}}(0)=1$, pero esta superposici\'on se
  encuentra m\'as alineada con una rama polarit\'onica. La generaci\'on de ramas
  fot\'on-fon\'on distinguibles es m\'as gradual y la pureza decrece con mayor
  lentitud, aunque finalmente se aproxima a $1/2$. El hamiltoniano efectivo
  reproduce de manera pr\'acticamente id\'entica la din\'amica del modelo HC
  completo.}
\end{figure}

\section{Entrelazamiento tripartito y eigenestados GHZ}

El sistema tiene tres grados de libertad. Los espacios de la cavidad y de la
vibraci\'on son infinitodimensionales, por lo que estrictamente no son qubits.
Sin embargo, dentro de un doblete fijo $(n,m)$ solo aparecen dos estados
relevantes de cada grado de libertad:
\begin{equation}
  \mathcal H_{\mathrm{el}}
  =
  \operatorname{span}\{\ket e,\ket g\},
\end{equation}
\begin{equation}
  \mathcal H_{\mathrm{cav}}^{(n)}
  =
  \operatorname{span}\{\ket{n-1},\ket n\},
\end{equation}
\begin{equation}
  \mathcal H_{\mathrm{vib}}^{(m)}
  =
  \operatorname{span}\{\ket{m-1},\ket m\}.
\end{equation}
Por esta raz\'on, el subespacio efectivo puede tratarse como un sistema de tres
qubits codificados:
\begin{equation}
  \tilde{\mathcal H}_{\mathrm{ef}}^{(n,m)}
  \subset
  \mathbb C_{\mathrm{el}}^2
  \otimes
  \mathbb C_{\mathrm{cav}}^2
  \otimes
  \mathbb C_{\mathrm{vib}}^2.
\end{equation}
El qubit electr\'onico es un qubit f\'isico de dos niveles. Los otros dos son
qubits codificados en pares de estados de Fock consecutivos.

Los eigenestados del hamiltoniano efectivo son
\begin{equation}
  \ket{E_{nm}^{+}}
  =
  \mathcal C_{nm}\ket{+,n,m-1}
  -
  \mathcal S_{nm}\ket{-,n,m},
\end{equation}
\begin{equation}
  \ket{E_{nm}^{-}}
  =
  \mathcal S_{nm}\ket{+,n,m-1}
  +
  \mathcal C_{nm}\ket{-,n,m},
\end{equation}
donde
\begin{equation}
  \mathcal C_{nm}
  =
  \sqrt{
  \frac{\Omega_{nm}^{\mathrm{ef}}+\delta_n}
  {2\Omega_{nm}^{\mathrm{ef}}}
  },
  \qquad
  \mathcal S_{nm}
  =
  \sqrt{
  \frac{\Omega_{nm}^{\mathrm{ef}}-\delta_n}
  {2\Omega_{nm}^{\mathrm{ef}}}
  }.
\end{equation}

En resonancia polarit\'on-fon\'on,
\begin{equation}
  \delta_n=0
  \quad\Longrightarrow\quad
  \mathcal C_{nm}=\mathcal S_{nm}=\frac1{\sqrt2}.
\end{equation}
Los eigenestados se reducen a
\begin{equation}
  \ket{E_{nm}^{+}}
  =
  \frac{
  \ket{+,n,m-1}-\ket{-,n,m}
  }{\sqrt2},
\end{equation}
\begin{equation}
  \ket{E_{nm}^{-}}
  =
  \frac{
  \ket{+,n,m-1}+\ket{-,n,m}
  }{\sqrt2}.
\end{equation}

Regresando a la base electr\'onica, fot\'onica y fon\'onica,
\begin{equation}
\begin{aligned}
  \ket{E_{nm}^{+}}
  ={}&
  \frac12
  \Big(
  \ket{e,n-1,m-1}
  +
  \ket{g,n,m-1}
  \\
  &\hspace{2.1cm}
  -
  \ket{e,n-1,m}
  +
  \ket{g,n,m}
  \Big)
  \\
  ={}&
  \ket e\ket{n-1}
  \frac{\ket{m-1}-\ket m}{2}
  +
  \ket g\ket n
  \frac{\ket{m-1}+\ket m}{2}.
\end{aligned}
\end{equation}
De manera an\'aloga,
\begin{equation}
\begin{aligned}
  \ket{E_{nm}^{-}}
  ={}&
  \frac12
  \Big(
  \ket{e,n-1,m-1}
  +
  \ket{g,n,m-1}
  \\
  &\hspace{2.1cm}
  +
  \ket{e,n-1,m}
  -
  \ket{g,n,m}
  \Big)
  \\
  ={}&
  \ket e\ket{n-1}
  \frac{\ket{m-1}+\ket m}{2}
  +
  \ket g\ket n
  \frac{\ket{m-1}-\ket m}{2}.
\end{aligned}
\end{equation}

Definimos los estados vibracionales
\begin{equation}
  \ket{\phi_m^{\pm}}
  =
  \frac{\ket{m-1}\pm\ket m}{\sqrt2}.
\end{equation}
Entonces
\begin{equation}
  \boxed{
  \ket{E_{nm}^{+}}
  =
  \frac1{\sqrt2}
  \left(
  \ket e\ket{n-1}\ket{\phi_m^{-}}
  +
  \ket g\ket n\ket{\phi_m^{+}}
  \right).
  }
\end{equation}
El segundo eigenestado es
\begin{equation}
  \boxed{
  \ket{E_{nm}^{-}}
  =
  \frac1{\sqrt2}
  \left(
  \ket e\ket{n-1}\ket{\phi_m^{+}}
  +
  \ket g\ket n\ket{\phi_m^{-}}
  \right).
  }
\end{equation}

Estos estados son de tipo Greenberger-Horne-Zeilinger. Cada t\'ermino de la
superposici\'on cambia simult\'aneamente el estado electr\'onico, el n\'umero de
fotones y el estado vibracional codificado. Para verificar que las dos ramas
son ortogonales en cada subsistema se usa
\begin{equation}
  \braket{e|g}=0,
  \qquad
  \braket{n-1|n}=0,
\end{equation}
y
\begin{equation}
\begin{aligned}
  \braket{\phi_m^{-}|\phi_m^{+}}
  &=
  \frac12
  (\bra{m-1}-\bra m)
  (\ket{m-1}+\ket m)
  \\
  &=
  \frac12(1+0-0-1)
  =0.
\end{aligned}
\end{equation}

Introduciendo los estados l\'ogicos
\begin{equation}
  \ket0_{\mathrm{el}}=\ket e,
  \qquad
  \ket1_{\mathrm{el}}=\ket g,
\end{equation}
\begin{equation}
  \ket0_{\mathrm{cav}}=\ket{n-1},
  \qquad
  \ket1_{\mathrm{cav}}=\ket n,
\end{equation}
\begin{equation}
  \ket0_{\mathrm{vib}}=\ket{\phi_m^{-}},
  \qquad
  \ket1_{\mathrm{vib}}=\ket{\phi_m^{+}},
\end{equation}
se obtiene
\begin{equation}
  \boxed{
  \ket{E_{nm}^{+}}
  =
  \frac{\ket{000}+\ket{111}}{\sqrt2}.
  }
\end{equation}
El otro eigenestado es
\begin{equation}
  \ket{E_{nm}^{-}}
  =
  \frac{\ket{001}+\ket{110}}{\sqrt2}
\end{equation}
con esta elecci\'on de estados l\'ogicos. Es localmente equivalente al GHZ
can\'onico mediante una compuerta $\hat X$ aplicada al qubit vibracional.

\section{Pureza reducida de los estados GHZ}

Para medir el entrelazamiento tripartito se considera cada subsistema frente a
los otros dos:
\begin{equation}
  \mathrm{vib}\,|\,(\mathrm{el,cav}),
  \qquad
  \mathrm{cav}\,|\,(\mathrm{el,vib}),
  \qquad
  \mathrm{el}\,|\,(\mathrm{cav,vib}).
\end{equation}
Esto no significa que cada pareja est\'e entrelazada por s\'i misma. Significa
que cada subsistema individual est\'a m\'aximamente entrelazado con los otros
dos considerados como un \'unico sistema compuesto.

Para $\ket{E_{nm}^{+}}$, la matriz de densidad total es
\begin{equation}
\begin{aligned}
  \hat\rho^{+}
  =\frac12\Big(&
  \ket{e,n-1,\phi_m^{-}}\bra{e,n-1,\phi_m^{-}}
  \\
  &+
  \ket{e,n-1,\phi_m^{-}}\bra{g,n,\phi_m^{+}}
  \\
  &+
  \ket{g,n,\phi_m^{+}}\bra{e,n-1,\phi_m^{-}}
  \\
  &+
  \ket{g,n,\phi_m^{+}}\bra{g,n,\phi_m^{+}}
  \Big).
\end{aligned}
\end{equation}
Al trazar sobre el estado electr\'onico, los t\'erminos cruzados desaparecen
por $\braket{e|g}=0$. La matriz reducida cavidad-vibraci\'on es
\begin{equation}
  \hat\rho_{\mathrm{cav-vib}}
  =
  \frac12
  \left(
  \ket{n-1,\phi_m^{-}}\bra{n-1,\phi_m^{-}}
  +
  \ket{n,\phi_m^{+}}\bra{n,\phi_m^{+}}
  \right).
\end{equation}
Los dos proyectores son ortogonales. Por ello,
\begin{equation}
  \hat\rho_{\mathrm{cav-vib}}^2
  =
  \frac14
  \left(
  \ket{n-1,\phi_m^{-}}\bra{n-1,\phi_m^{-}}
  +
  \ket{n,\phi_m^{+}}\bra{n,\phi_m^{+}}
  \right),
\end{equation}
y
\begin{equation}
  \boxed{
  P(\hat\rho_{\mathrm{cav-vib}})
  =
  \operatorname{Tr}(\hat\rho_{\mathrm{cav-vib}}^2)
  =
  \frac12.
  }
\end{equation}
Trazar sobre cualquiera de los otros subsistemas produce el mismo valor de
pureza. Cada qubit codificado tiene una matriz reducida m\'aximamente mezclada.
Las matrices de dos qubits obtenidas al descartar el tercero son, en cambio,
mezclas cl\'asicas de dos productos ortogonales y no contienen
entrelazamiento bipartito. El entrelazamiento reside genuinamente en la
partici\'on tripartita.

\section{Generaci\'on din\'amica de estados GHZ}

Consideremos un polarit\'on que evoluciona bajo el hamiltoniano efectivo del
doblete $(n,m)$:
\begin{equation}
  \hat H_{\mathrm{ef}}^{(n,m)}
  =
  \hbar A_{nm}\Id
  +
  \frac{\hbar}{2}
  \left(
  \delta_n\hat\sigma_z
  -
  \kappa_m\hat\sigma_x
  \right),
\end{equation}
en la base
\begin{equation}
  \left\{
  \ket{+,n,m-1},
  \ket{-,n,m}
  \right\}.
\end{equation}
Aqu\'i
\begin{equation}
  \kappa_m=\nu\lambda\sqrt m,
  \qquad
  \Omega_{nm}^{\mathrm{ef}}
  =
  \sqrt{\delta_n^2+\kappa_m^2}.
\end{equation}

El operador de evoluci\'on es
\begin{equation}
  \hat U_{nm}(t)
  =
  \exp\left(
  -\frac{\ii}{\hbar}
  \hat H_{\mathrm{ef}}^{(n,m)}t
  \right).
\end{equation}
El t\'ermino proporcional a la identidad conmuta con el resto, de modo que la
exponencial se factoriza:
\begin{equation}
  \hat U_{nm}(t)
  =
  e^{-\ii A_{nm}t}
  \exp\left[
  -\frac{\ii}{2}
  \left(
  \delta_n\hat\sigma_z
  -
  \kappa_m\hat\sigma_x
  \right)t
  \right].
\end{equation}

Definimos
\begin{equation}
  \hat M_{nm}
  =
  \delta_n\hat\sigma_z
  -
  \kappa_m\hat\sigma_x.
\end{equation}
Como las matrices de Pauli satisfacen
\begin{equation}
  \hat\sigma_z^2=\hat\sigma_x^2=\Id,
  \qquad
  \{\hat\sigma_z,\hat\sigma_x\}=0,
\end{equation}
se tiene
\begin{equation}
\begin{aligned}
  \hat M_{nm}^2
  &=
  \delta_n^2\hat\sigma_z^2
  +
  \kappa_m^2\hat\sigma_x^2
  -
  \delta_n\kappa_m
  \left(
  \hat\sigma_z\hat\sigma_x
  +
  \hat\sigma_x\hat\sigma_z
  \right)
  \\
  &=
  \left(
  \delta_n^2+\kappa_m^2
  \right)\Id
  \\
  &=
  \left(
  \Omega_{nm}^{\mathrm{ef}}
  \right)^2\Id.
\end{aligned}
\end{equation}
En consecuencia, las potencias pares son
\begin{equation}
  \hat M_{nm}^{2q}
  =
  \left(
  \Omega_{nm}^{\mathrm{ef}}
  \right)^{2q}\Id,
\end{equation}
y las impares son
\begin{equation}
  \hat M_{nm}^{2q+1}
  =
  \left(
  \Omega_{nm}^{\mathrm{ef}}
  \right)^{2q}\hat M_{nm}.
\end{equation}

La exponencial de una matriz se define mediante su serie de potencias:
\begin{equation}
  e^{\hat X}
  =
  \sum_{r=0}^{\infty}
  \frac{\hat X^r}{r!}.
\end{equation}
Separando las potencias pares de $\hat M_{nm}$,
\begin{equation}
\begin{aligned}
  \sum_{q=0}^{\infty}
  \frac{1}{(2q)!}
  \left(-\frac{\ii t}{2}\right)^{2q}
  \hat M_{nm}^{2q}
  &=
  \sum_{q=0}^{\infty}
  \frac{(-1)^q}{(2q)!}
  \left(
  \frac{\Omega_{nm}^{\mathrm{ef}}t}{2}
  \right)^{2q}
  \Id
  \\
  &=
  \cos\left(
  \frac{\Omega_{nm}^{\mathrm{ef}}t}{2}
  \right)
  \Id.
\end{aligned}
\end{equation}
Para las potencias impares,
\begin{equation}
\begin{aligned}
  &\sum_{q=0}^{\infty}
  \frac{1}{(2q+1)!}
  \left(-\frac{\ii t}{2}\right)^{2q+1}
  \hat M_{nm}^{2q+1}
  \\
  &\qquad=
  -\ii
  \frac{\hat M_{nm}}{\Omega_{nm}^{\mathrm{ef}}}
  \sum_{q=0}^{\infty}
  \frac{(-1)^q}{(2q+1)!}
  \left(
  \frac{\Omega_{nm}^{\mathrm{ef}}t}{2}
  \right)^{2q+1}
  \\
  &\qquad=
  -\ii
  \frac{\hat M_{nm}}{\Omega_{nm}^{\mathrm{ef}}}
  \sin\left(
  \frac{\Omega_{nm}^{\mathrm{ef}}t}{2}
  \right).
\end{aligned}
\end{equation}
Por tanto, el operador de evoluci\'on completo es
\begin{equation}
  \boxed{
  \hat U_{nm}(t)
  =
  e^{-\ii A_{nm}t}
  \left[
  \cos\left(
  \frac{\Omega_{nm}^{\mathrm{ef}}t}{2}
  \right)\Id
  -
  \ii
  \frac{
  \delta_n\hat\sigma_z-
  \kappa_m\hat\sigma_x
  }{\Omega_{nm}^{\mathrm{ef}}}
  \sin\left(
  \frac{\Omega_{nm}^{\mathrm{ef}}t}{2}
  \right)
  \right].
  }
\end{equation}

Partimos del estado inicial
\begin{equation}
  \ket{\tilde\Psi(0)}
  =
  \ket{+,n,m-1}.
\end{equation}
En la base del doblete,
\begin{equation}
  \hat\sigma_z\ket{+,n,m-1}
  =
  \ket{+,n,m-1},
\end{equation}
\begin{equation}
  \hat\sigma_x\ket{+,n,m-1}
  =
  \ket{-,n,m}.
\end{equation}
Al aplicar $\hat U_{nm}(t)$ se obtiene
\begin{equation}
\begin{aligned}
  \ket{\tilde\Psi(t)}
  =e^{-\ii A_{nm}t}\Bigg\{&
  \left[
  \cos\left(
  \frac{\Omega_{nm}^{\mathrm{ef}}t}{2}
  \right)
  -
  \ii
  \frac{\delta_n}{\Omega_{nm}^{\mathrm{ef}}}
  \sin\left(
  \frac{\Omega_{nm}^{\mathrm{ef}}t}{2}
  \right)
  \right]
  \ket{+,n,m-1}
  \\
  &+
  \ii
  \frac{\kappa_m}{\Omega_{nm}^{\mathrm{ef}}}
  \sin\left(
  \frac{\Omega_{nm}^{\mathrm{ef}}t}{2}
  \right)
  \ket{-,n,m}
  \Bigg\}.
\end{aligned}
\end{equation}

Omitiendo la fase global, definimos
\begin{equation}
  a(t)
  =
  \cos\left(
  \frac{\Omega_{nm}^{\mathrm{ef}}t}{2}
  \right)
  -
  \ii
  \frac{\delta_n}{\Omega_{nm}^{\mathrm{ef}}}
  \sin\left(
  \frac{\Omega_{nm}^{\mathrm{ef}}t}{2}
  \right),
\end{equation}
\begin{equation}
  b(t)
  =
  \ii
  \frac{\kappa_m}{\Omega_{nm}^{\mathrm{ef}}}
  \sin\left(
  \frac{\Omega_{nm}^{\mathrm{ef}}t}{2}
  \right).
\end{equation}
As\'i,
\begin{equation}
  \ket{\tilde\Psi(t)}
  =
  a(t)\ket{+,n,m-1}
  +
  b(t)\ket{-,n,m}.
\end{equation}

Regresando a la base electr\'onica desnuda,
\begin{equation}
\begin{aligned}
  \ket{\tilde\Psi(t)}
  =
  \frac1{\sqrt2}\Big[&
  \ket e\ket{n-1}
  \left(
  a(t)\ket{m-1}+b(t)\ket m
  \right)
  \\
  &+
  \ket g\ket n
  \left(
  a(t)\ket{m-1}-b(t)\ket m
  \right)
  \Big].
\end{aligned}
\end{equation}
Los estados vibracionales que acompa\~nan a las ramas electr\'onicas son
\begin{equation}
  \ket{\xi_e(t)}
  =
  a(t)\ket{m-1}+b(t)\ket m,
\end{equation}
\begin{equation}
  \ket{\xi_g(t)}
  =
  a(t)\ket{m-1}-b(t)\ket m.
\end{equation}
Su producto interno es
\begin{equation}
  \braket{\xi_g(t)|\xi_e(t)}
  =
  |a(t)|^2-|b(t)|^2.
\end{equation}
Para que el estado sea un GHZ con entrelazamiento m\'aximo se requiere que
estas ramas sean ortogonales:
\begin{equation}
  |a(t)|^2-|b(t)|^2=0.
\end{equation}
Como la evoluci\'on es unitaria,
\begin{equation}
  |a(t)|^2+|b(t)|^2=1.
\end{equation}
Combinando ambas ecuaciones se obtiene la condici\'on
\begin{equation}
  \boxed{
  |a(t)|^2=|b(t)|^2=\frac12.
  }
\end{equation}

En resonancia,
\begin{equation}
  \delta_n=0
  \quad\Longrightarrow\quad
  \Omega_{nm}^{\mathrm{ef}}=\kappa_m.
\end{equation}
El estado evolucionado se simplifica a
\begin{equation}
  \ket{\tilde\Psi(t)}
  =
  e^{-\ii A_{nm}t}
  \left[
  \cos\left(
  \frac{\kappa_mt}{2}
  \right)
  \ket{+,n,m-1}
  +
  \ii
  \sin\left(
  \frac{\kappa_mt}{2}
  \right)
  \ket{-,n,m}
  \right].
\end{equation}
Las probabilidades de las dos ramas polarit\'onicas son
\begin{equation}
  P_+(t)
  =
  \cos^2\left(
  \frac{\kappa_mt}{2}
  \right),
  \qquad
  P_-(t)
  =
  \sin^2\left(
  \frac{\kappa_mt}{2}
  \right).
\end{equation}
Para generar un GHZ m\'aximo se requiere
\begin{equation}
  \cos^2\left(
  \frac{\kappa_mt}{2}
  \right)
  =
  \sin^2\left(
  \frac{\kappa_mt}{2}
  \right)
  =
  \frac12.
\end{equation}
Esto ocurre cuando
\begin{equation}
  \frac{\kappa_mt}{2}
  =
  \frac{(2q+1)\pi}{4},
  \qquad
  q=0,1,2,\ldots
\end{equation}
Por tanto, los tiempos de generaci\'on son
\begin{equation}
  \boxed{
  t_q^{\mathrm{GHZ}}
  =
  \frac{(2q+1)\pi}{2\kappa_m}
  =
  \frac{(2q+1)\pi}{2\nu\lambda\sqrt m}.
  }
\end{equation}
Es decir, partiendo de un polarit\'on puro aparece un estado GHZ en cada uno
de los tiempos $t_q^{\mathrm{GHZ}}$.

Para el primer tiempo, y omitiendo la fase global, el estado en la base
desnuda es
\begin{equation}
  \boxed{
  \begin{aligned}
  \ket{\tilde\Psi(t_0^{\mathrm{GHZ}})}
  =
  \frac1{\sqrt2}\Bigg[&
  \ket e\ket{n-1}
  \frac{\ket{m-1}+\ii\ket m}{\sqrt2}
  \\
  &+
  \ket g\ket n
  \frac{\ket{m-1}-\ii\ket m}{\sqrt2}
  \Bigg].
  \end{aligned}
  }
\end{equation}
Los estados vibracionales de las dos ramas son ortogonales, al igual que los
estados electr\'onicos y fot\'onicos. El estado es, por tanto, un GHZ can\'onico
salvo fases locales.

\end{document}
